\chapter{Hoja de ruta de implementación}
\label{cap:hoja-ruta}

Este capítulo propone un orden de trabajo. Está pensado para alguien que tiene un
modelo estático funcionando y un semestre por delante. Cada etapa termina con un
criterio verificable: si no se cumple, no conviene pasar a la siguiente.

\section{Etapa 0: el estado estacionario del modelo que ya tienes}

\paragraph{Qué hacer} Escribir en código el modelo estático que ya existe, pero
estructurado como el estado estacionario de un modelo dinámico: una función que,
dados los parámetros y los precios, devuelva todas las cantidades agregadas.

\paragraph{Criterio de salida} El código reproduce exactamente los números del
modelo estático escrito en papel, y las identidades contables (el balance de cada
banco, la identidad $X-G=\mu$, la suma de masas igual a uno) se cumplen a
precisión de máquina.

\paragraph{Cuánto toma} Una o dos semanas. Es tiempo bien gastado: el 80\% de los
errores que aparecerían después se atrapan aquí.

\section{Etapa 1: cerrar las decisiones pendientes}

\paragraph{Qué hacer} Convertir en decisiones todo lo que ahora es parámetro
impuesto. En el modelo de la parte V, eso significa derivar la liquidez
voluntaria de una condición de primer orden en lugar de fijarla. Para cada
decisión que se cierre, hay que escribir explícitamente el problema de
optimización, su condición de primer orden, y verificar numéricamente que la
solución la satisface.

\paragraph{Criterio de salida} Ninguna variable de elección es un parámetro
libre, y para cada una se puede decir qué fuerzas económicas la determinan y en
qué dirección responde a cada precio.

\section{Etapa 2: añadir el estado endógeno individual}

\paragraph{Qué hacer} Introducir el estado que acumula --- el patrimonio, en el
caso bancario --- con su ley de movimiento y su problema dinámico. Resolver el
estado estacionario completo: iteración hacia atrás, distribución estacionaria,
agregación.

\paragraph{Criterio de salida}
\begin{enumerate}[leftmargin=1.4em,itemsep=1pt]
\item La distribución converge a $10^{-14}$ y suma uno.
\item La masa en el último punto de la grilla es despreciable (digamos, menor
que $10^{-8}$). Si no, la grilla está atando y hay que extenderla o revisar la
condición de crecimiento.
\item El modelo reproduce la versión sin estado endógeno cuando se apaga el
mecanismo que lo hace relevante.
\item Las magnitudes agregadas son plausibles: ratio de capital, rentabilidad,
apalancamiento.
\end{enumerate}

\paragraph{Cuánto toma} Dos a cuatro semanas. Es la etapa con más trabajo
conceptual nuevo.

\section{Etapa 3: los jacobianos del bloque}

\paragraph{Qué hacer} Implementar el algoritmo fake news para el bloque, y
verificarlo.

\paragraph{Criterio de salida} \emph{Este es el criterio innegociable del
proyecto.} El jacobiano calculado con el algoritmo rápido coincide con el del
método directo, a $T=30$ o $40$, con error relativo menor que $10^{-6}$ en
\emph{todos} los pares insumo--producto. Además, todas las identidades
analíticas que el modelo implique se verifican.

\begin{trampa}
La tentación de saltarse la verificación contra el método directo es enorme,
porque el algoritmo rápido ``parece funcionar'': produce matrices con la forma
correcta y signos razonables. Casi todos los errores de implementación de este
método --- el orden invertido de los operadores, un índice desfasado en la
recursión, una distribución mal alineada --- producen jacobianos que se ven bien
y están mal. Escribir el método directo cuesta veinte líneas y es la única
defensa real.
\end{trampa}

\section{Etapa 4: el equilibrio general}

\paragraph{Qué hacer} Armar el grafo, elegir incógnitas y objetivos, resolver.

\paragraph{Criterio de salida}
\begin{enumerate}[leftmargin=1.4em,itemsep=1pt]
\item El número de condición de $\mathbf{H}_U$ es razonable.
\item Las respuestas en impacto de las variables predeterminadas son cero.
\item La condición de vaciado omitida por la ley de Walras da residuo numérico
nulo.
\item La solución no cambia al pasar de $T$ a $1.5T$.
\end{enumerate}

\section{Etapa 5: disciplinar la calibración}

Aquí es donde el proyecto deja de ser un ejercicio y empieza a decir algo. La
pregunta que organiza esta etapa es: \emph{¿qué dato identifica a qué parámetro?}

La tabla \ref{tab:identificacion} es un ejemplo de cómo organizarla para el
modelo bancario de la parte V. El punto no son las entradas concretas sino el
ejercicio: para cada parámetro, un momento, y para cada momento, una razón por la
que es informativo sobre ese parámetro y no sobre otro.

\begin{table}[htbp]
\centering\small\setlength{\tabcolsep}{4pt}
\caption{Qué dato disciplina a qué parámetro en el modelo bancario.}
\label{tab:identificacion}
\begin{tabular}{@{}p{0.19\textwidth}p{0.33\textwidth}p{0.40\textwidth}@{}}
\toprule
\textbf{Parámetro} & \textbf{Momento} & \textbf{Por qué es informativo}\\
\midrule
$a_e$ (exposición) & volatilidad del crecimiento de depósitos, por banco &
es directamente la dispersión del choque de retiro, en el corte transversal\\
$\tau_e$ (acceso) & saldos interbancarios sobre activos, por banco &
un banco con mejor acceso opera más en el mercado\\
$w_e$ (dolarización) & participación de depósitos en ME, por banco &
observable directamente\\
$\lambda$ (apalancamiento) & ratio de capital del sistema & contable\\
$\eta$, $p$ (estado agregado) & frecuencia y magnitud de episodios de salida de
depósitos en ME & serie agregada\\
$v_F$, $\bar\iota^u$ (respaldo) & uso de las operaciones de liquidez del banco
central en episodios de tensión & es la variable que el modelo llama $C^u$\\
$\beta$, $\varsigma$ (banquero) & rentabilidad y tasa de salida de bancos &
fijan el crecimiento y la escala\\
\bottomrule
\end{tabular}
\end{table}

\begin{idea}
Nótese que casi todos los momentos de la tabla son \emph{transversales}, no de
series de tiempo. Eso no es casualidad: la heterogeneidad se identifica con
variación entre unidades. Una estrategia de calibración que solo use agregados no
puede separar parámetros que afectan a los agregados de la misma forma, como se
advirtió en el capítulo \ref{cap:estimacion}.
\end{idea}

\section{Etapa 6: los ejercicios que valen la pena}

Una vez que el modelo funciona, estos son los ejercicios que un modelo dinámico
puede hacer y uno estático no. Están ordenados por relación entre valor y
esfuerzo.

\begin{enumerate}[leftmargin=1.4em]
\item \textbf{Respuesta a un cambio de encaje, anticipado y no anticipado.}
Comparar la columna $s=0$ con la columna $s=8$ del jacobiano de equilibrio
general. Esto responde directamente a una pregunta de diseño de política: ¿sirve
anunciar con antelación?
\item \textbf{Incidencia dinámica por tipo.} Cómo se reparte el ajuste entre
clases de banco, y cómo cambia ese reparto a lo largo de la transición.
\item \textbf{Erosión del capital.} Cuánto del efecto de largo plazo de una
política viene del canal de acumulación y cuánto del efecto contemporáneo. Se
mide apagando el estado endógeno.
\item \textbf{Transición entre regímenes de encaje.} Usar el método del capítulo
\ref{cap:nolineal} para calcular la trayectoria completa de una reforma
permanente.
\item \textbf{A qué tamaño de choque se rompe la linealidad.} Encontrar el
choque de escasez que agota la capacidad de respaldo, y mostrar cuánto se separa
la respuesta no lineal de la lineal en ese punto.
\end{enumerate}

\section{Qué hacer si el tiempo se acaba}

Un plan realista tiene que contemplar que no todo entre. Si hay que recortar, el
orden de prioridad sugerido es:

\begin{itemize}[leftmargin=1.4em]
\item \textbf{Conservar siempre:} la etapa 3 completa, con su verificación. Un
jacobiano verificado es un resultado defendible aunque no haya equilibrio general.
\item \textbf{Recortar primero:} la estimación. Calibrar es perfectamente
defendible en una tesis de pregrado, y la estimación bayesiana agrega mucho
riesgo por poco contenido adicional.
\item \textbf{Recortar segundo:} las transiciones no lineales. Son valiosas pero
son un complemento.
\item \textbf{Simplificar antes de abandonar:} si el modelo completo con dos
monedas y tres tipos no converge a tiempo, un modelo con una moneda y dos tipos
que \emph{funciona y está verificado} vale infinitamente más que uno completo que
no corre.
\end{itemize}
